\section*{Exercises}
Exercises~\ref{ex:13.1}--\ref{ex:13.4} are core exercises, and Exercises~\ref{ex:13.5} and~\ref{ex:13.6} are transfer exercises. Complete solutions are in Appendix~\ref{app:solc}, under ``Week~13.''
\begin{exercise}\label{ex:13.1}A random variable takes the values $0$, $2$, and $5$ with probabilities $1/2$, $1/3$, and $1/6$. Compute its expectation, and name a value of the variable that is at least as large as the expectation.\end{exercise}
\begin{exercise}\label{ex:13.2}Label each vertex of a triangle with $0$ or $1$, and take the cut whose left side consists of the vertices labeled $0$. List the possible sizes of this cut, and compute the average cut size over all eight labelings.\end{exercise}
\begin{exercise}\label{ex:13.3}Run the construction of Section~\ref{ch13:sec:construction} on the graph with vertices $1,2,3,4$ and edges $12$, $23$, $34$, $41$ (here $12$ stands for the edge $\{1,2\}$, and so on). Place the vertices in the order $1,2,3,4$, and break ties by choosing the left side. At each step, record $\ell$, $r$, the side chosen, and the new value of $F$. How large is the final cut, and how does it compare with the guarantee of Theorem~\ref{thm:cut}?\end{exercise}
\begin{exercise}\label{ex:13.4}In a finite probability space, four bad events each have probability at most $0.2$. Prove that some outcome of positive probability avoids all four. Then suppose instead that each bad event has probability at most $0.25$. Does the same conclusion follow? Explain, using examples.\end{exercise}
\begin{exercise}\label{ex:13.5}Give an example of a random variable whose expectation is at least $5$ but which takes the value $0$ at some outcome of positive probability. This shows that a large expectation does not force every value to be large. Then state what the expectation does guarantee about the values of a random variable.\end{exercise}
\begin{exercise}\label{ex:13.6}Let $X$ be a nonnegative random variable with $\E X=3$. Use Markov's inequality to bound $\Prob(X\ge12)$. Then give an example of such a random variable, on a finite probability space, for which $\Prob(X\ge12)$ equals your bound.\end{exercise}
